% section: 発展パターン

\section{発展パターン}
ここでは更に発展的なパターンについて紹介する.
扱われる漸化式は,一部の難関大において時々出題されるようなものである.
基本的に誘導がついた状態で出題されるため,ひと目見て解けるようになる必要はそこまで高くないが,知っていて損はない.
\subsection{ガウス型}
\subsubsection{下準備}
まず,ガウス記号について説明しておこう.
ガウス記号とは,\,$\lfloor x \rfloor$または,\,$[x]$で表される記号で,これを床関数という.
意味はある実数$x$について,以下を満たす$k$のことを指す.
\[
  k \le x < k+1
\]
日本語で言えば実数$x$の整数部分,すなわち切り捨てである.
また,類似の記号として,\,$\lceil x \rceil$が存在する.これは,切り上げを意味し,天井関数と呼ばれている.
きちんと数式で書くと
\[
  k-1 < x \le k
\]
を満たす$k$である.

ガウス記号を漸化式に含む数列について考える前に,一般項にガウスを含む数列について考える.
例えばこのような数列を考えよう.
\[
  a_1=1,\quad a_{n+1}=\lfloor a_n+\sqrt{3} \rfloor
\]
一見難しそうな見た目をしているが,その実この漸化式は
\[
  a_1=1,\quad a_{n+1}=a_n+2
\]
と同値である.
なぜならば,\,$a_n$は常に整数を取り,そこに$\sqrt{3}$を足してから小数を切り捨てている.
結局2を足し続けているだけなのだ.

では少し漸化式から離れて数列にガウス記号が含まれている場合を考える.
\[
  a_{n}=\lfloor \sqrt{n} \rfloor
\]
この数列はどのような数列になるだろうか.
結果から言えばこの一般項からなる数列は群数列を成す.
証明は以下のとおりである.
\begin{quote}
  \noindent \textbf{【命題1】}
  \begin{quote}
    一般項が
    \[
      a_n=\lfloor \sqrt{n} \rfloor\quad(n\ge1)
    \]
    と定義される数列$a_n$は群数列か.
  \end{quote}

  \noindent \textbf{【証明1】}
  \begin{quote}
    数列$a_n$に対し,
    \[
      a_n=k
    \]
    を満たす正整数$k$が存在するときの$n$が取りうる範囲について考える.床関数の定義から,
    \[
      k\le\sqrt{n}<k+1
    \]
    $n,k\ge1$から,
    \begin{align*}
       & k^2\le n <(k+1)^2                      \\
       & \therefore\quad k^2\le n \le (k+1)^2-1
    \end{align*}
    これを満たす$n$の範囲は
    \[
      \{(k+1)^2-1\}-k^2+1=2k+1
    \]
    このことから,ある$a_n=k$となるときの$n$は$2k+1$項分だけ存在する.
    従って,\,$k$に依存して連続して同じ数を持つ項の数が決定されることから,一般項が
    で定義される数列$a_n$は群数列であることが示された. \qed
  \end{quote}
\end{quote}
実際に実験をして検証してみよう.
初項から順に書き出してみると,
\[
  1,1,1,2,2,2,2,2,3,3,3,3,3,3,3,4
\]
となっており,理論と一致している.

もう一例ぐらい考えてみよう.
\[
  a_{n}=\left\lfloor \frac{n^2+2n+1}{4} \right\rfloor -\left\lfloor \frac{n^2-2n+1}{4} \right\rfloor
\]
複雑な見た目をした一般項だ.こういう時は実験として初項から書き出してみるといい.\,$a_1$から順に,
\[
  1,2,3,4,5,\cdots
\]
不思議なことに,この数列の一般項は$a_n=n$の一般項と$a_5$までは一致する.
更に進めてみても同様に一致し続ける.
その真偽を検証しよう.
\begin{quote}
  \noindent \textbf{【命題2】}
  \begin{quote}
    一般項が
    \[
      a_{n}=\left\lfloor \frac{n^2+2n+1}{4} \right\rfloor -\left\lfloor \frac{n^2-2n+1}{4} \right\rfloor
    \]
    と定義される数列$a_n$は初項1,公差1の等差数列と同値か.
  \end{quote}

  \noindent \textbf{【証明2】}
  \begin{quote}
    任意の実数$x$と任意の整数$k$について,
    \[
      \left\lfloor x+k \right\rfloor=\left\lfloor x \right\rfloor +k
    \]
    が成り立つ.
    このことから,
    \[
      X=\frac{(n-1)^2}{4}
    \]
    とおくと,与式は
    \begin{align*}
      a_n & =\left\lfloor X+\frac{4n}{4} \right\rfloor-\left\lfloor X \right\rfloor \\
          & =\left\lfloor X \right\rfloor+n-\left\lfloor X \right\rfloor            \\
          & =n
    \end{align*}
    従って題意は示された.\qed
  \end{quote}
\end{quote}
$n=4k$と置いて$\mod$を使うような解法も考えられる.

いずれにせよ,ガウスを含む漸化式は細かな（つまりは小数の）変化を均すことがこれらの議論からわかる.
\subsubsection{ガウス型の解説}
ガウス記号を含む数列を通して,ガウス記号の扱いにある程度慣れたところで本題に入る.

この漸化式の一般項はどのように表されるだろうか.
\[
  a_{n+1}=a_n+\lfloor \sqrt{n} \rfloor
\]
この漸化式は階差数列型だから,
\[
  a_n=a_1+\sum_{k=1}^{n-1} \lfloor \sqrt{k} \rfloor
\]
$b_n=\lfloor \sqrt{n} \rfloor$を考える.この数列は証明2より群数列を成す.
第$m$群の総和は,
\[
  \sum_{i=1}^{2m+1} m=2m^2+m
\]
従って第$l$群までの総和は,
\[
  \sum_{j=1}^{l}  2j^2+j
\]
第$n$項が第何群に属するかは,
\subsection{隣接$n$項間漸化式型}
隣接2項間や隣接3項間の漸化式を今まで扱った.
これをより広い範囲に適用してみよう.

\subsection{虚数解型隣接3項間漸化式}
今までの隣接3項間漸化式は,特性方程式の解が実数解のみを有した.
ここでは虚数解となるような場合を考える.
\subsection{非隣接2項間漸化式型}
\subsection{複数組み合わせ型}
